% Auto-generated from raw.txt
% Usage:
%   \vocabentry{word}{/ipa/ (pos)}{definition}{example}{note}

\vocabentry{Infrastructure}
{/ˈɪnfrəstrʌktʃər/ (n)}
{The basic physical and organizational structures needed for the operation of a society.}
{Developing world-class transport infrastructure is essential for attracting foreign investment.}
{Đây là một danh từ dùng để chỉ cơ sở hạ tầng như cầu đường, hệ thống điện nước.}

\vocabentry{Congestion}
{/kənˈdʒestʃən/ (n)}
{The state of being extremely full or blocked, especially with traffic.}
{New policies such as congestion charges aim to reduce the number of vehicles entering the city center.}
{Đây là một danh từ dùng để chỉ sự tắc nghẽn hoặc ùn tắc giao thông.}

\vocabentry{Commute}
{/kəˈmjuːt/ (v/n)}
{To travel some distance between one's home and place of work on a regular basis.}
{Many people prefer living in the suburbs despite having to endure a long daily commute to the office.}
{Đây là một động từ/danh từ chỉ việc đi làm hàng ngày bằng phương tiện giao thông.}

\vocabentry{Public transport}
{/ˌpʌblɪk ˈtrænspɔːrt/ (n)}
{Buses, trains, and other forms of transport that are available to the public, charge set fares, and run on fixed routes.}
{An efficient public transport system is essential for reducing carbon emissions in urban areas.}
{Đây là một cụm danh từ chỉ phương tiện giao thông công cộng.}

\vocabentry{Logistics}
{/ləˈdʒɪstɪks/ (n)}
{The detailed coordination of a complex operation involving many people, facilities, or supplies.}
{Effective logistics planning is essential for reducing waste and costs in the supply chain.}
{Đây là một danh từ dùng để chỉ ngành hậu cần/vận tải.}

\vocabentry{Connectivity}
{/ˌkɑːnekˈtɪvəti/ (n)}
{The state or quality of being connected or connective.}
{The new high-speed rail link will significantly improve connectivity between the capital and the northern provinces.}
{Trong giao thông, đây là danh từ chỉ khả năng kết nối giữa các vùng miền hoặc các loại hình vận tải.}

\vocabentry{Accessibility}
{/əkˌsesəˈbɪləti/ (n)}
{The quality of being able to be reached or entered.}
{Improving the accessibility of the subway system for disabled passengers is a top priority.}
{Đây là một danh từ dùng để chỉ tính dễ tiếp cận của một địa điểm hoặc phương tiện giao thông.}

\vocabentry{Intermodal}
{/ˌɪntərˈmoʊdl/ (adj)}
{Involving two or more different modes of transport in conveying goods.}
{Intermodal shipping containers can be easily transferred from ships to trains and then to trucks.}
{Đây là một tính từ dùng để chỉ vận tải đa phương thức, kết hợp nhiều loại phương tiện.}

\vocabentry{Gridlock}
{/ˈɡrɪdlɑːk/ (n)}
{A traffic jam in which no vehicle can move in any direction.}
{A minor accident during rush hour caused total gridlock throughout the downtown area.}
{Đây là một danh từ dùng để chỉ tình trạng giao thông tê liệt hoàn toàn.}

\vocabentry{Modernization}
{/ˌmɑːrdənəˈzeɪʃn/ (n)}
{The process of adapting something to modern needs or habits.}
{The modernization of the air traffic control system will enhance safety and reduce flight delays.}
{Đây là một danh từ dùng để chỉ sự hiện đại hóa hạ tầng hoặc hệ thống vận hành.}

\vocabentry{Pedestrianization}
{/pəˌdestriənaɪˈzeɪʃn/ (n)}
{The process of making a street or area for use by people walking only.}
{The pedestrianization of the historic district has encouraged more tourists and local shoppers.}
{Đây là một danh từ dùng để chỉ việc chuyển đổi thành phố đi bộ, cấm phương tiện cơ giới.}

\vocabentry{Traffic-free}
{/ˈtræfɪk friː/ (adj)}
{Restricted to pedestrians and sometimes cyclists; no motor vehicles allowed.}
{The town center is a traffic-free zone on weekends to promote outdoor activities.}
{Đây là một tính từ dùng để chỉ khu vực không có phương tiện giao thông cơ giới.}

\vocabentry{Bottleneck}
{/ˈbɑːtlnek/ (n)}
{A situation that causes delay in a process or system.}
{The narrow bridge acts as a bottleneck, causing long tailbacks during the morning rush hour.}
{Trong giao thông, đây là danh từ chỉ điểm nghẽn, nơi đường hẹp gây ùn tắc.}

\vocabentry{Bypass}
{/ˈbaɪpæs/ (n/v)}
{A road that goes around a town or congested area; to go around or avoid something.}
{Construction of the northern bypass will begin early next year.}
{Đây là một danh từ/động từ dùng để chỉ đường tránh hoặc tránh đi vòng.}

\vocabentry{Junction}
{/ˈdʒʌŋkʃn/ (n)}
{A place where roads or railway lines meet, cross, or diverge.}
{The accident happened at the busy junction of the two main highways.}
{Đây là một danh từ dùng để chỉ nút giao.}

\vocabentry{Viaduct}
{/ˈvaɪədʌkt/ (n)}
{A long bridge carrying a road or railway across a valley, river, or other obstacle.}
{The modern railway viaduct is a striking feature of the local landscape.}
{Đây là một danh từ dùng để chỉ cầu cạn.}

\vocabentry{Flyover}
{/ˈflaɪoʊvər/ (n)}
{A bridge that carries one road over another.}
{The flyover was closed for several hours after a truck hit the protective barrier.}
{Đây là một danh từ dùng để chỉ cầu vượt.}

\vocabentry{Subway}
{/ˈsʌbweɪ/ (n)}
{An underground railway system, especially in a city.}
{Londoners often refer to their subway system as the "Tube."}
{Đây là một danh từ dùng để chỉ tàu điện ngầm.}

\vocabentry{Light rail}
{/laɪt reɪl/ (n)}
{An urban rail transport system using relatively light trains, often with street-level sections.}
{Light rail systems are often seen as a greener alternative to diesel buses.}
{Đây là một danh từ dùng để chỉ tàu điện nhẹ.}

\vocabentry{Freight}
{/freɪt/ (n)}
{Goods transported in bulk by truck, train, ship, or aircraft.}
{Air freight is much more expensive than sea freight but also much faster.}
{Đây là một danh từ dùng để chỉ hàng hóa vận chuyển.}

\vocabentry{Fleet}
{/fliːt/ (n)}
{A group of ships or a collection of vehicles operated by a company or organization.}
{The government plans to replace its entire fleet of official cars with electric models by 2030.}
{Đây là một danh từ dùng để chỉ hạm đội/đội xe.}

\vocabentry{Decentralization}
{/ˌdiːˌsentrələˈzeɪʃn/ (n)}
{The transfer of control of an activity or organization to several local offices or authorities.}
{Decentralization of government offices can help reduce the pressure on urban transport networks.}
{Trong hạ tầng, đây là việc dời các dịch vụ ra xa trung tâm.}

\vocabentry{Redevelopment}
{/ˌriːdɪˈveləpmənt/ (n)}
{The action or process of developing something again or differently.}
{The redevelopment of the old dockyards has transformed the area into a vibrant commercial hub.}
{Đây là một danh từ dùng để chỉ việc quy hoạch lại hoặc xây dựng lại một khu vực hạ tầng cũ.}

\vocabentry{Throughput}
{/ˈθruːpʊt/ (n)}
{The amount of material or items passing through a system or process.}
{The port has increased its annual container throughput by 15\% through automation.}
{Trong vận tải, đây là lưu lượng hàng hóa hoặc hành khách thông qua cảng/nhà ga.}

\vocabentry{Transit}
{/ˈtrænzɪt/ (n/v)}
{The carrying of people or things from one place to another.}
{Light rail transit is a cost-effective alternative to expensive subway systems.}
{Đây là một danh từ dùng để chỉ sự vận chuyển tài nguyên.}

\vocabentry{Autonomous}
{/ɔːˈtɑːnəməs/ (adj)}
{(Of a vehicle) having the power of self-government; self-driving.}
{Autonomous buses are currently being tested in several smart cities around the world.}
{Đây là một tính từ dùng để chỉ các phương tiện tự hành, không người lái.}

\vocabentry{Electrification}
{/ɪˌlektrɪfɪˈkeɪʃn/ (n)}
{The action of process of charging something with electricity.}
{The electrification of the main rail line has resulted in faster and more reliable train services.}
{Đây là một danh từ dùng để chỉ quá trình điện khí hóa, ví dụ hệ thống đường sắt.}

\vocabentry{Interchange}
{/ˈɪntərtʃeɪndʒ/ (n)}
{A complex junction where several roads or motorways meet.}
{The bus interchange is located next to the main railway station for easy passenger transfer.}
{Đây là một danh từ dùng để chỉ nút giao thông lập thể, nơi nhiều con đường gặp nhau.}

\vocabentry{Maintenance}
{/ˈmeɪntənəns/ (n)}
{The process of maintaining or preserving someone or something.}
{High maintenance costs are a major challenge for aging subway systems.}
{Đây là một danh từ dùng để chỉ sự bảo trì thiết bị năng lượng.}

\vocabentry{Capacity}
{/kəˈpæsəti/ (n)}
{The maximum amount that something can contain.}
{During peak hours, the subway system often operates at full capacity.}
{Trong giao thông, đây là danh từ chỉ sức chứa hoặc khả năng đáp ứng lưu lượng.}

\vocabentry{Vulnerability}
{/ˌvʌlnərəˈbɪləti/ (n)}
{The quality or state of being exposed to the possibility of being attacked or harmed.}
{The flood exposed the vulnerability of the coastal railway line to rising sea levels.}
{Trong hạ tầng, đây là sự dễ bị tổn thương trước thiên tai hoặc sự cố.}

\vocabentry{Resilience}
{/rɪˈzɪliəns/ (n)}
{The capacity to recover quickly from difficulties; toughness.}
{Building climate resilience into transport infrastructure is a key challenge for engineers today.}
{Trong hạ tầng, đây là khả năng chống chịu và phục hồi nhanh sau sự cố.}

\vocabentry{Standardization}
{/ˌstændərdəˈzeɪʃn/ (n)}
{The process of making things of the same type all have the same basic features.}
{Standardization of container sizes has made international shipping much more efficient.}
{Đây là một danh từ dùng để chỉ sự áp dụng tiêu chuẩn chung cho các thiết bị và hệ thống năng lượng}

\vocabentry{Compliance}
{/kəmˈplaɪəns/ (n)}
{The action or fact of complying with a wish or command.}
{Monitoring the compliance of transport companies with labor laws is essential.}
{Đây là một danh từ dùng để chỉ sự tuân thủ các quy định pháp luật.}

\vocabentry{Sustainability}
{/səˌsteɪnəˈbɪləti/ (n)}
{The ability to be maintained at a certain rate or level.}
{The long-term sustainability of urban growth depends on a reliable transport infrastructure.}
{Đây là một danh từ dùng để chỉ tính bền vững của đô thị về mặt môi trường và kinh tế.}

\vocabentry{Feasibility}
{/ˌfiːzəˈbɪləti/ (n)}
{The state or degree of being easily or conveniently done.}
{A detailed feasibility study was conducted before the construction of the underwater tunnel began.}
{Đây là một danh từ dùng để chỉ tính khả thi của một dự án xây dựng hạ tầng.}

\vocabentry{Utilization}
{/ˌjuːtələˈzeɪʃn/ (n)}
{The action of making practical and effective use of something.}
{Improving the utilization of off-peak transport can help to balance the load on the network.}
{Đây là một danh từ dùng để chỉ sự tận dụng tối đa nguồn lực năng lượng hiện có}

\vocabentry{Optimization}
{/ˌɑːptɪməˈzeɪʃn/ (n)}
{The action of making the best or most effective use of a situation or resource.}
{Optimization of maintenance schedules can help to reduce down-time for transport equipment.}
{Đây là một danh từ dùng để chỉ quá trình tối ưu hóa hiệu suất hệ thống năng lượng}

\vocabentry{Standard}
{/ˈstændərd/ (n/adj)}
{A level of quality or attainment.}
{All new vehicles must meet strict international standards for exhaust emissions.}
{Đây là một danh từ/tính từ dùng để chỉ các tiêu chuẩn kỹ thuật hoặc an toàn.}

\vocabentry{Regulation}
{/ˌreɡjuˈleɪʃn/ (n)}
{A rule or directive made and maintained by an authority.}
{New regulations have been introduced to control the noise from night-time flights.}
{Đây là một danh từ dùng để chỉ quy định hoặc sự điều tiết của chính phủ.}

\vocabentry{Implementation}
{/ˌɪmplɪmenˈteɪʃn/ (n)}
{The process of putting a decision or plan into effect; execution.}
{The implementation of smart road signs has improved traffic flow during peak hours.}
{Đây là một danh từ dùng để chỉ sự thực hiện chính sách hoặc dự án năng lượng}

\vocabentry{Significant}
{/sɪɡˈnɪfɪkənt/ (adj)}
{Sufficiently great or important to be worthy of attention.}
{There has been a significant shift towards the use of bicycles in major European cities.}
{Đây là một tính từ dùng để chỉ sự đáng kể hoặc quan trọng trong thay đổi hạ tầng.}

\vocabentry{Essential}
{/ɪˈsenʃl/ (adj)}
{Absolutely necessary; extremely important.}
{Reliable transport links are essential for the economic development of rural areas.}
{Đây là tính từ dùng để chỉ các dịch vụ hoặc hạ tầng thiết yếu.}

\vocabentry{Vital}
{/ˈvaɪtl/ (adj)}
{Absolutely necessary or important.}
{Public investment is vital for maintaining a safe and modern road network.}
{Đây là một tính từ dùng để chỉ tầm quan trọng sống còn.}

\vocabentry{Fundamental}
{/ˌfʌndəˈmentl/ (adj)}
{Forming a necessary base or core.}
{Clean energy is a fundamental component of any sustainable transport strategy.}
{Đây là một tính từ dùng để chỉ những yếu tố nền tảng của hạ tầng.}

\vocabentry{Drastic}
{/ˈdræstɪk/ (adj)}
{Likely to have a strong or far-reaching effect; radical and extreme.}
{Drastic measures are needed to tackle the growing problem of urban air pollution.}
{Đây là một tính từ dùng để chỉ sự thay đổi mạnh mẽ, quyết liệt trong giao thông.}

\vocabentry{Gradual}
{/ˈɡrædʒuəl/ (adj)}
{Taking place or progressing slowly or by degrees.}
{There has been a gradual increase in the number of electric cars on the road over the last few years.}
{Đây là tính từ dùng để chỉ sự thay đổi dần dần.}

\vocabentry{Rapid}
{/ˈræpɪd/ (adj)}
{Happening in a short time or at a great rate.}
{The rapid expansion of the highway network has fueled economic growth in the region.}
{Đây là một tính từ dùng để chỉ sự phát triển nhanh chóng của hạ tầng.}

\vocabentry{Widespread}
{/ˌwaɪdˈspred/ (adj)}
{Found or distributed over a large area or number of people.}
{There is widespread concern about the safety of self-driving vehicles on public roads.}
{Đây là một tính từ dùng để chỉ sự phổ biến rộng rãi.}

\vocabentry{Unprecedented}
{/ʌnˈpresɪdentɪd/ (adj)}
{Never done or known before.}
{The city is facing an unprecedented level of traffic congestion due to the metro construction.}
{Đây là một tính từ dùng để chỉ những thay đổi hoặc sự cố chưa từng có tiền lệ.}

\vocabentry{Paradigm}
{/ˈpærədaɪm/ (n)}
{A typical example or pattern of something; a model.}
{The shift from private cars to shared mobility represents a new paradigm in urban transport.}
{Trong giao thông, đây là danh từ dùng để chỉ một mô hình hoặc chuẩn mực phát triển.}

\vocabentry{Hub}
{/hʌb/ (n)}
{The effective center of an activity, region, or network.}
{The city is a major logistics hub, with excellent connections to the road and rail networks.}
{Đây là một danh từ dùng để chỉ trung tâm hoặc đầu mối toàn cầu.}

\vocabentry{Vessel}
{/ˈvesl/ (n)}
{A container in which liquid may be kept.}
{The coast guard is monitoring the movement of all commercial vessels in the area.}
{Trong ẩm thực, đây là danh từ chỉ dụng cụ đựng thực phẩm hoặc bình chứa.}

\vocabentry{Tarmac}
{/ˈtɑːrmæk/ (n)}
{A paved surface, especially an area at an airport where aircraft park or move.}
{The heat of the sun had softened the tarmac on the old country road.}
{Đây là một danh từ dùng để chỉ mặt đường nhựa/khu sân đỗ máy bay.}

\vocabentry{Pavement}
{/ˈpeɪvmənt/ (n)}
{A raised paved or asphalted path for pedestrians at the side of a road.}
{The city council is replacing the old concrete pavement with more attractive stone tiles.}
{Đây là một danh từ dùng để chỉ vỉa hè dành cho người đi bộ.}

\vocabentry{Carriage}
{/ˈkærɪdʒ/ (n)}
{A vehicle or compartment for carrying passengers, especially a railway carriage.}
{We managed to find two empty seats in the last carriage of the train.}
{Đây là một danh từ dùng để chỉ toa xe/hành khách hoặc khoang tàu.}

\vocabentry{Deck}
{/dek/ (n)}
{A floor-like platform on a ship; also a level of a bus or other vehicle.}
{The car deck of the ferry was full of tourists heading to the islands.}
{Đây là một danh từ dùng để chỉ boong tàu/tầng của phương tiện.}

\vocabentry{Chassis}
{/ˈʃæsi/ (n)}
{The frame of a vehicle, onto which the body and other parts are attached.}
{The truck's chassis was damaged in the accident and had to be completely replaced.}
{Đây là một danh từ dùng để chỉ khung gầm.}

\vocabentry{Suspension}
{/səˈspenʃn/ (n)}
{The system of springs and shock absorbers that connects a vehicle to its wheels.}
{High-quality suspension is essential for vehicles that carry fragile goods.}
{Đây là một danh từ dùng để chỉ hệ thống treo/giảm xóc.}

\vocabentry{Transmission}
{/trænzˈmɪʃn/ (n)}
{The action or process of transmitting something.}
{Manual transmission is still common in many parts of Europe and Asia.}
{Trong năng lượng là sự truyền tải điện năng đi xa.}

\vocabentry{Aeronautics}
{/ˌerəˈnɔːtɪks/ (n)}
{The science or practice of travel through the air.}
{He is studying aeronautics at university with the dream of designing faster planes.}
{Đây là một danh từ dùng để chỉ ngành hàng không.}

\vocabentry{Navigation}
{/ˌnævɪˈɡeɪʃn/ (n)}
{The process or activity of accurately ascertaining one's position and planning and following a route.}
{Good navigation skills are essential for sailors and pilots alike.}
{Hàng hải lịch sử.}

\vocabentry{Avionics}
{/ˌeɪviˈɑːnɪks/ (n)}
{Electronic equipment used in aircraft and spacecraft.}
{Modern avionics systems are incredibly complex and highly automated.}
{Đây là một danh từ dùng để chỉ hệ thống điện tử hàng không vũ trụ.}

\vocabentry{Altitude}
{/ˈæltɪtuːd/ (n)}
{The height of an object or point in relation to sea level or ground level.}
{The aircraft reached its cruising altitude of thirty-five thousand feet twenty minutes after takeoff.}
{Đây là một danh từ dùng để chỉ độ cao so với mặt nước biển hoặc bề mặt hành tinh.}

\vocabentry{Velocity}
{/vəˈlɑːsəti/ (n)}
{The speed of something in a given direction.}
{The rocket must reach a specific velocity to enter Earth's orbit.}
{Vận tốc.}

\vocabentry{Acceleration}
{/əkˌseləˈreɪʃn/ (n)}
{The increase in speed or rate of change of velocity.}
{The car's acceleration was so powerful that it reached one hundred kilometers per hour in five seconds.}
{Đây là một danh từ dùng để chỉ sự tăng tốc.}

\vocabentry{Braking}
{/ˈbreɪkɪŋ/ (n)}
{The action of slowing down or stopping a vehicle using brakes.}
{Anti-lock braking systems (ABS) have significantly improved road safety.}
{Đây là một danh từ dùng để chỉ việc phanh/giảm tốc.}

\vocabentry{Aerodynamics}
{/ˌerodaɪˈnæmɪks/ (n)}
{The study of how air moves and how it affects objects moving through it, especially aircraft.}
{Formula One teams spend millions of dollars every year on improving the aerodynamics of their cars.}
{Đây là một danh từ dùng để chỉ ngành khí động học.}

\vocabentry{Combustion}
{/kəmˈbʌstʃən/ (n)}
{The process of burning something.}
{Lean-burn combustion technology can reduce both fuel consumption and harmful emissions.}
{Đây là một danh từ dùng để chỉ sự đốt cháy nhiên liệu.}

\vocabentry{Hybrid}
{/ˈhaɪbrɪd/ (n/adj)}
{The offspring of two plants or animals of different species or varieties.}
{Hybrid buses are becoming more common in cities that want to improve air quality.}
{Đây là một danh từ dùng để chỉ cá thể lai.}

\vocabentry{Emissions}
{/iˈmɪʃnz/ (n)}
{The production and discharge of something, especially gas or radiation.}
{Zero-emissions vehicles are essential for meeting the city's environmental goals.}
{Đây là một danh từ dùng để chỉ khí thải, đặc biệt là khí gây ô nhiễm từ nhà máy hoặc xe cộ.}

\vocabentry{Particulate}
{/pɑːrˈtɪkjələt/ (n/adj)}
{Tiny particles, especially those suspended in air; relating to particles.}
{Filtering particulate matter from engine exhaust is a major technical challenge.}
{Đây là một danh từ/tính từ dùng để chỉ hạt bụi mịn hoặc dạng hạt.}

\vocabentry{Fumes}
{/fjuːmz/ (n)}
{Strong-smelling or dangerous gas or smoke.}
{Exhaust fumes can cause irritation to the eyes, nose, and throat.}
{Đây là một danh từ dùng để chỉ khói độc.}

\vocabentry{Pipeline}
{/ˈpaɪplaɪn/ (n)}
{A long pipe, typically underground, for conveying oil, gas, etc. over long distances.}
{The construction of the new pipeline was delayed by environmental protests.}
{Đây là một danh từ dùng để chỉ đường ống dẫn năng lượng.}

\vocabentry{Quarry}
{/ˈkwɔːri/ (n/v)}
{A place, typically a large, deep pit, from which stone or other materials are or have been extracted.}
{Noise from the local quarry has been a cause of complaints among nearby residents.}
{Đây là một danh từ dùng để chỉ mỏ đá.}

\vocabentry{Refining}
{/rɪˈfaɪnɪŋ/ (n)}
{The process of removing impurities or unwanted elements from a substance.}
{Refining capacity in the region has been increased to meet the growing demand for fuel.}
{Trong năng lượng là quá trình lọc dầu hoặc tinh chế quặng.}

\vocabentry{Extraction}
{/ɪkˈstrækʃn/ (n)}
{The action of taking out something, especially using effort or force.}
{The extraction of natural gas from shale has revolutionized the energy market.}
{Đây là một danh từ dùng để chỉ sự khai thác tài nguyên như dầu mỏ hoặc khoáng sản.}

\vocabentry{Subsidize}
{/ˈsʌbsɪdaɪz/ (v)}
{Support (an organization or activity) financially.}
{In many countries, the government subsidizes student travel to make education more accessible.}
{Động từ của subsidy, chỉ việc trợ cấp.}

\vocabentry{Efficiency}
{/ɪˈfɪʃnsi/ (n)}
{The state or quality of being efficient.}
{Fuel efficiency has become a key selling point for new car manufacturers.}
{Đây là một danh từ dùng để chỉ hiệu quả kinh tế hoặc năng suất.}

\vocabentry{Terminal}
{/ˈtɜːrmɪnl/ (n/adj)}
{A point where something ends or is distributed.}
{The airport terminal was crowded with holidaymakers heading to the sun.}
{Đây là một danh từ dùng để chỉ cảng hoặc trạm trung chuyển tài nguyên.}

\vocabentry{Concourse}
{/ˈkɑːŋɔːrs/ (n)}
{A large open area inside or in front of a public building.}
{Shops and restaurants line the main concourse of the international rail station.}
{Đây là một danh từ dùng để chỉ sảnh lớn ở sân bay hoặc nhà ga.}

\vocabentry{Platform}
{/ˈplætfɔːrm/ (n)}
{A standard for a computer system's hardware and software.}
{The train for London will depart from platform number three in five minutes.}
{Đây là một danh từ dùng để chỉ nền tảng công nghệ.}

\vocabentry{Bay}
{/beɪ/ (n)}
{A broad inlet of the sea where the land curves inward.}
{Trucks must wait in the designated loading bay until their goods are ready.}
{Đây là một danh từ dùng để chỉ vịnh.}

\vocabentry{Siding}
{/ˈsaɪdɪŋ/ (n)}
{A short track beside the main railway line used for passing or storing trains.}
{Several old locomotives were kept on a siding near the old railway museum.}
{Đây là một danh từ dùng để chỉ đường ray nhánh.}

\vocabentry{Sleepers}
{/ˈsliːpərz/ (n)}
{Beams (wooden or concrete) laid across a railway track to support the rails.}
{Workers were busy replacing the old wooden sleepers with new concrete ones.}
{Đây là một danh từ dùng để chỉ tà vẹt đường sắt.}

\vocabentry{Ballast}
{/ˈbæləst/ (n)}
{Heavy material used to provide stability, especially gravel under railway tracks or weight in a ship.}
{Fresh ballast was spread along the new section of the railway track.}
{Đây là một danh từ dùng để chỉ vật dằn/đá ballast giúp ổn định.}

\vocabentry{Gauge}
{/ɡeɪdʒ/ (n/v)}
{A measuring device; also the distance between rails on a railway; to measure or estimate.}
{Standard gauge is used by most major railway systems around the world.}
{Đây là một danh từ/động từ dùng để chỉ thước đo/khổ ray hoặc đo lường.}

\vocabentry{Rolling stock}
{/ˈroʊlɪŋ stɑːk/ (n)}
{Railway vehicles, such as locomotives, carriages, and wagons.}
{Modernizing the rolling stock is a key part of the railway improvement plan.}
{Đây là một danh từ dùng để chỉ đoàn toa xe/thiết bị phương tiện đường sắt.}

\vocabentry{Locomotive}
{/ˌloʊkəˈmoʊtɪv/ (n/adj)}
{The engine that pulls a train; relating to movement or propulsion.}
{The historic steam locomotive is now used for special tourist trips on the weekend.}
{Đây là một danh từ/tính từ dùng để chỉ đầu máy xe lửa hoặc thuộc về chuyển động.}

\vocabentry{Wagon}
{/ˈwæɡən/ (n)}
{A railway vehicle used for transporting goods; a large vehicle for carrying goods.}
{The delivery wagon was pulled by a pair of strong horses through the narrow streets.}
{Đây là một danh từ dùng để chỉ toa hàng/xe chở hàng.}

\vocabentry{Dredging}
{/ˈdredʒɪŋ/ (n)}
{The process of removing mud, sand, or sediment from the bottom of a river, harbor, or channel to deepen it.}
{Dredging operations are ongoing to deepen the entrance to the main harbor.}
{Đây là một danh từ dùng để chỉ hoạt động nạo vét.}

\vocabentry{Quay}
{/kiː/ (n)}
{A stone or concrete platform beside water where ships load and unload.}
{Passengers boarded the ferry at the main quay near the city center.}
{Đây là một danh từ dùng để chỉ cầu cảng/bờ kè bốc dỡ.}

\vocabentry{Lock}
{/lɑːk/ (n/v)}
{A gated section of a canal used to raise or lower boats between different water levels; to secure something.}
{It took nearly an hour for the large barge to pass through the canal lock.}
{Đây là một danh từ/động từ dùng để chỉ âu tàu/khóa.}

\vocabentry{Estuary}
{/ˈestʃueri/ (n)}
{The tidal mouth of a large river where it meets the stream.}
{The estuary is a vital habitat for many species of fish and birds.}
{Đây là một danh từ dùng để chỉ cửa sông.}

\vocabentry{Canal}
{/kəˈnæl/ (n)}
{An artificial waterway built to allow boats or ships to travel inland.}
{Many historic canals are now popular destinations for leisure boating and walking.}
{Đây là một danh từ dùng để chỉ kênh đào/kênh nhân tạo.}

\vocabentry{Harbor}
{/ˈhɑːrbər/ (n/v)}
{A place on the coast where vessels may find shelter.}
{Several small islands are located just outside the natural harbor.}
{Đây là một danh từ dùng để chỉ bến cảng.}

\vocabentry{Pier}
{/pɪr/ (n)}
{A platform on pillars projecting from the shore into the water, used for docking or walking.}
{Fishing boats are moored along the side of the stone pier.}
{Đây là một danh từ dùng để chỉ cầu tàu/bến tàu.}

\vocabentry{Breakwater}
{/ˈbreɪkwaːtər/ (n)}
{A barrier built to protect a shore or harbor from the force of waves.}
{The city is building a new breakwater to protect coastal homes from erosion.}
{Đây là một danh từ dùng để chỉ đê chắn sóng.}

\vocabentry{Berth}
{/bɜːrθ/ (n/v)}
{A ship's assigned place at a dock; also a place to sleep on a ship or train.}
{The ship reached its designated berth at exactly three o'clock in the afternoon.}
{Đây là một danh từ/động từ dùng để chỉ chỗ neo đậu hoặc đưa tàu vào bến.}

\vocabentry{Wharf}
{/wɔːrf/ (n)}
{A structure built along the shore where ships can dock to load and unload.}
{Large cranes on the wharf were busy loading grain onto the ship.}
{Đây là một danh từ dùng để chỉ cầu cảng/bến cảng.}

\vocabentry{Jetty}
{/ˈdʒeti/ (n)}
{A structure extending into the water to protect a harbor or provide a place for boats to dock.}
{We watched the sunset from the end of the stone jetty.}
{Đây là một danh từ dùng để chỉ cầu cảng/đê chắn nhỏ vươn ra biển.}

\vocabentry{Hangar}
{/ˈhæŋər/ (n)}
{A large building where aircraft are stored and maintained.}
{The aircraft was towed from the runway back to its hangar for maintenance.}
{Đây là một danh từ dùng để chỉ nhà chứa máy bay.}

\vocabentry{Runway}
{/ˈrʌnweɪ/ (n)}
{A strip of land at an airport on which aircraft take off and land.}
{High-intensity lights are used to illuminate the runway during night landings.}
{Đây là một danh từ dùng để chỉ đường băng.}

\vocabentry{Taxiway}
{/ˈtæksiweɪ/ (n)}
{A path on an airport along which aircraft taxi between the runway and the terminal.}
{The ground crew guided the plane along the taxiway to the departure gate.}
{Đây là một danh từ dùng để chỉ đường lăn.}

\vocabentry{Control tower}
{/kənˈtroʊl ˌtaʊər/ (n)}
{A tall building at an airport from which air traffic controllers monitor and guide aircraft.}
{From the top of the control tower, the controllers have a clear view of the entire airport.}
{Đây là một danh từ dùng để chỉ tháp điều khiển không lưu.}

\vocabentry{Fuselage}
{/ˈfjuːsəlɑːʒ/ (n)}
{The main body of an aircraft, to which the wings and tail are attached.}
{The fuselage was heavily damaged when the plane made an emergency landing.}
{Đây là một danh từ dùng để chỉ thân máy bay.}

\vocabentry{Wing}
{/wɪŋ/ (n)}
{The flat structure on an aircraft that produces lift and enables it to fly.}
{I prefer to sit near the wing of the plane because it feels more stable.}
{Đây là một danh từ dùng để chỉ cánh máy bay.}

\vocabentry{Cockpit}
{/ˈkɑːkpɪt/ (n)}
{The compartment in an aircraft where the pilot sits and controls the plane.}
{Unauthorized personnel are strictly prohibited from entering the cockpit during the flight.}
{Đây là một danh từ dùng để chỉ buồng lái.}

